\section{Collective Eigenmodes and Retained Outlook Formulae}\label{app:collective}
The conditional formulas and balanced-release definition from former Sections~38--43 retain their original numbers.

\subsection{Collective eigenmodes}\label{sec:modes}
For an admitted global linear update $U$, eigenvectors or invariant subspaces define collective eigenmodes. A reversible quantum representation has unit-modulus eigenvalues. A reduced dissipative representation can exhibit contraction when complementary registers are omitted. Particle-like interpretation additionally requires preparation, transport, stability, interaction, and readout correspondences.

\subsection{Thermal occupation}\label{sec:temperature}
Given assigned energies $E_j$, maximizing $-\sum_jp_j\log p_j$ at fixed normalization and mean energy yields
% Original equation number(s) 49; expression unchanged.
\setcounter{equation}{48}
\begin{equation}
 p_j=\frac{e^{-\beta E_j}}{Z(\beta)},
 \qquad Z(\beta)=\sum_j e^{-\beta E_j}.
\end{equation}

The identification $\beta=(k_B\Theta)^{-1}$ requires the stated energy scale and thermal interpretation. The variational calculation does not supply equilibration dynamics or make arbitrary occupation profiles thermal.

\subsection{Balanced release}\label{sec:release}
\begingroup\renewcommand{\thetheorem}{40.1}
\setcounter{theorem}{0}
\begin{definition}[Balanced release]
Balanced release is efficient local conversion of internal excitation into admissible propagating field modes, with the complementary relational information required for the event preserved in the complete accounting.
\end{definition}

\endgroup
Here the original phrase ``field modes'' refers to collective eigenmodes, not to Mode. For a normalized preparation $\psi$, an admitted linear operation $U$, and an outgoing projector $P_{\rm out}$, the retained illustration is
% Original equation number(s) 50; expression unchanged.
\setcounter{equation}{49}
\begin{equation}
\eta(\psi)=\|P_{\mathrm{out}}U\psi\|^2.
\end{equation}

It is an outgoing weight. An energy-efficiency interpretation requires an energy-compatible representation and normalization. A large expectation of a separate excitation operator does not force this weight to be large.

